\chapter{Relatividad Especial}

\section{Transformación de Galileo}

\begin{observacion}[Concepto de relatividad]
El objetivo de cualquier teoría de la relatividad, sea Galileana-
Newtoniana o Einsteiniana, es describir la manera en que diferentes
observadores ven un mismo fenómeno. La pregunta natural es: ¿son las
leyes de la naturaleza las mismas para todos los observadores?, o
¿existen estados de movimiento privilegiados o sistemas de referencia
privilegiados?
\end{observacion}

La teoría de Newton postula la existencia de sistemas de referencia
privilegiados, como puede extraerse de la primera ley del movimiento:

\begin{quote}
\emph{``Todo cuerpo continúa en su estado de reposo o de movimiento
uniforme y rectilíneo a menos que este sea obligado a cambiar dicho
estado por fuerzas que actúan sobre él.''}
\end{quote}

Esto significa que existe un conjunto de cuerpos privilegiados, aquellos
sobre los que no actúan fuerzas. Una vez encontrado un \textbf{sistema de
referencia inercial} (S.R.I.), se tiene que todo otro sistema que esté
en reposo o se mueva con velocidad constante con respecto de dicho
S.R.I. es también un S.R.I.

\begin{definicion}[Transformación de Galileo]
Es una transformación de coordenadas que conecta un sistema inercial con
otro sistema inercial. Consideremos los sistemas de referencia $K$ y
$K'$, donde $K'$ se mueve con respecto de $K$ con velocidad constante
$V$ a lo largo del eje $x$.
\end{definicion}

\begin{center}
\begin{tikzpicture}[scale=1]
  \draw[-Stealth] (0,0) -- (2.2,0) node[right] {$x$};
  \draw[-Stealth] (0,0) -- (0.9,1.6) node[above] {$y$};
  \draw[-Stealth] (0,0) -- (-1.1,0.7) node[left] {$z$};
  \node at (-0.3,-0.3) {$K$};
  \draw[-Stealth] (4,0) -- (6.2,0) node[right] {$x'$};
  \draw[-Stealth] (4,0) -- (4.9,1.6) node[above] {$y'$};
  \draw[-Stealth] (4,0) -- (2.9,0.7) node[left] {$z'$};
  \node at (3.7,-0.3) {$K'$};
  \draw[-Stealth,thick,colorej] (2.4,-0.8) -- (3.6,-0.8) node[midway,below] {$V$};
  \draw[fill] (1.3,1.1) circle (1.5pt) node[above right] {$P$};
  \node at (1.7,-1.3) {\footnotesize$(x,y,z,t)$};
  \node at (5.7,-1.3) {\footnotesize$(x',y',z',t')$};
\end{tikzpicture}
\end{center}

\begin{definicion}
Las medidas del observador en $K$ y las del observador en $K'$ están
relacionadas por medio de la transformación
\begin{equation}
t' = t, \qquad x' = x+Vt, \qquad y'=y, \qquad z'=z, \tag{1}
\end{equation}
donde hemos supuesto que los observadores usan las mismas unidades de
medida, relojes idénticos y perfectamente sincronizados.
\end{definicion}

La ecuación (1) es conocida como \textbf{transformación de Galileo}, la
que es la expresión matemática del \textbf{principio de relatividad de
Galileo}:

\begin{quote}
\emph{``Ningún experimento mecánico realizado en el interior (aislado)
de un sistema de Galileo permite poner de manifiesto el movimiento de
este sistema con relación a cualquier otro sistema de Galileo.''}
\end{quote}

\subsection{Invariancia de las ecuaciones de Newton}

Si $K$ y $K'$ son S.R.I., entonces es válido escribir
\[
m\vec a = \vec F \quad\text{(válido en $K$)}, \qquad m\vec a\,' = \vec
F\,' \quad\text{(válido en $K'$)}.
\]

\begin{proposicionc}[Transformación de velocidad y aceleración]
De (1) se obtiene
\[
v' = \frac{dx'}{dt} = \frac{d}{dt}(x+Vt) = \frac{dx}{dt}+V = v+V, \qquad
a' = \frac{dv'}{dt} = \frac{d}{dt}(v+V) = \frac{dv}{dt} = a.
\]
\end{proposicionc}

La velocidad claramente es relativa, o sea cambia de valor al pasar de
un S.R.I. a otro S.R.I., pero la \textbf{aceleración es invariante} bajo
el grupo de transformaciones de Galileo.

Dado que la masa es, en mecánica de Newton, un invariante, $m=m'$, se
sigue $\vec F=\vec F\,'$. Por lo tanto, \textbf{las ecuaciones de Newton
son invariantes bajo el grupo de transformaciones de Galileo}. Tanto en
$K$ como en $K'$ se cumplen exactamente las mismas leyes: no existe la
manera de discernir experimentalmente si es $K$ el que se traslada con
respecto de $K'$ o si es $K'$ el que se traslada con respecto a $K$.

\begin{observacion}[Simetría de la transformación]
Los siguientes casos son equivalentes:
\[
\begin{aligned}
&K' \text{ se mueve con } +V \text{ respecto de } K: \\
&\quad t'=t,\ x'=x+Vt,\ y'=y,\ z'=z,
\end{aligned}
\qquad\qquad
\begin{aligned}
&K \text{ se mueve con } -V \text{ respecto de } K': \\
&\quad t=t',\ x=x'-Vt,\ y=y',\ z=z'.
\end{aligned}
\]
\end{observacion}

En mecánica de Newton los conceptos de velocidad y trayectoria son
\textbf{relativos}: la velocidad depende del S.R.I. ($v'=v+V$) y la
trayectoria depende del S.R.I.

\begin{ejemplo}[Caída libre]
La caída libre de un objeto en $K'$ está descrita por la trayectoria
\[
x'=Vt', \qquad y'=y_0-\frac12 gt'^2, \qquad y=0, \qquad z=0,
\]
de donde, eliminando $t'$,
\[
y' = y_0 - \frac{g\,x'^2}{2V^2} \qquad \text{(parábola)}.
\]
Es decir, un objeto que cae verticalmente en $K$ describe una parábola
en $K'$.
\end{ejemplo}

Así tenemos que, en mecánica de Newton:
\begin{itemize}
\item La velocidad es relativa. \qquad La trayectoria es relativa.
\item La aceleración es absoluta. \qquad La masa es absoluta. \qquad La
fuerza es absoluta.
\item La distancia entre dos puntos es absoluta.
\item El intervalo de tiempo entre dos acontecimientos es absoluto, es
decir, el tiempo es absoluto.
\end{itemize}

\subsection{Las ecuaciones de Maxwell bajo el grupo de Galileo}

¿Satisfacen las leyes de la física el principio de relatividad? Hemos
visto que la mecánica de Newton es invariante bajo el grupo de Galileo.
Ahora nos preguntamos sobre el comportamiento de los sistemas
electromagnéticos bajo el grupo de Galileo.

Entre 1854 y 1865, J.\,C. Maxwell desarrolló la teoría electromagnética,
cuyas ecuaciones fundamentales son
\[
\text{(Gauss eléctrica)}\quad \vec\nabla\cdot\vec E = \frac{\rho}
{\varepsilon_0}, \qquad \text{(Faraday)}\quad \frac{\partial\vec
B}{\partial t}+\vec\nabla\times\vec E=0,
\]
\[
\text{(Ampère mod.)}\quad \vec\nabla\times\vec B - \varepsilon_0\mu_0
\frac{\partial\vec E}{\partial t} = \mu_0\vec J, \qquad
\text{(Gauss magnética)}\quad \vec\nabla\cdot\vec B=0.
\]

En el vacío ($\rho=0$, $\vec J=0$) se tiene
\[
\vec\nabla\cdot\vec E=0 \ \text{(2)}, \qquad \frac{\partial\vec
B}{\partial t}+\vec\nabla\times\vec E=0 \ \text{(3)}, \qquad
\vec\nabla\times\vec B-\varepsilon_0\mu_0\frac{\partial\vec E}{\partial
t}=0 \ \text{(4)}, \qquad \vec\nabla\cdot\vec B=0 \ \text{(5)}.
\]
Estas ecuaciones forman un conjunto compatible: (2) y (5) son
compatibles entre sí, y (3) y (4) son compatibles entre sí. Para obtener
ecuaciones que dependan solo de $\vec E$ (y solo de $\vec B$) subimos el
orden de las derivadas en uno.

\begin{proposicionc}
Tomando el rotacional a (4) y usando la identidad
$\vec\nabla\times\vec\nabla\times\vec A = \vec\nabla(\vec\nabla\cdot
\vec A) - \nabla^2\vec A$ junto con (3) y (5), se obtiene
\[
\nabla^2\vec B - \varepsilon_0\mu_0\, \frac{\partial^2\vec B}{\partial
t^2} = 0 . \tag{8}
\]
Análogamente, tomando el rotacional a (3) y usando (2) y (4), se
obtiene
\[
\nabla^2\vec E - \varepsilon_0\mu_0\, \frac{\partial^2\vec E}{\partial
t^2} = 0 . \tag{10}
\]
\end{proposicionc}

\begin{proof}
Tomando el rotacional a (4),
\[
\vec\nabla\times\vec\nabla\times\vec B = \varepsilon_0\mu_0\,
\vec\nabla\times\frac{\partial\vec E}{\partial t} = \varepsilon_0\mu_0\,
\frac{\partial}{\partial t}\big(\vec\nabla\times\vec E\big) =
\varepsilon_0\mu_0\, \frac{\partial}{\partial t}\left(-\frac{\partial
\vec B}{\partial t}\right) = -\varepsilon_0\mu_0\, \frac{\partial^2\vec
B}{\partial t^2}, \tag{6}
\]
donde se usó (3). Por otro lado, usando la identidad vectorial con
$\vec A=\vec B$ y (5),
\[
\vec\nabla\times\vec\nabla\times\vec B = \vec\nabla(\vec\nabla\cdot\vec
B) - \nabla^2\vec B = -\nabla^2\vec B . \tag{7}
\]
Igualando (6) y (7) se obtiene (8). El cálculo para $\vec E$, tomando el
rotacional a (3) y usando (4) y (2), es análogo y da (9) y luego (10).
\qedhere
\end{proof}

Ambas ecuaciones tienen la forma general
\[
\nabla^2\Psi - \varepsilon_0\mu_0\, \frac{\partial^2\Psi}{\partial t^2}
= 0, \tag{11}
\]
donde $\Psi$ representa cualquiera de las seis componentes $E_x, E_y,
E_z, B_x, B_y, B_z$.

\begin{proposicionc}[Rapidez de la luz]
Con $\varepsilon_0\approx 9\times10^9\ \text{N}\,\text{m}^2/\text{C}^2$
y $\mu_0=4\pi\times10^{-7}\ \text{N}/\text{A}^2$, se define
\[
c := \frac{1}{\sqrt{\varepsilon_0\mu_0}} = 299\,792\,458\ \text{m/s}
\approx 3\times10^8\ \text{m/s}, \tag{12}
\]
que es la rapidez de propagación de la luz en el vacío, de modo que
\[
\nabla^2\Psi - \frac{1}{c^2}\, \frac{\partial^2\Psi}{\partial t^2} = 0
. \tag{13}
\]
\end{proposicionc}

\begin{observacion}
En 1850, Weber y Kohlrausch determinaron el valor de
$1/\sqrt{\varepsilon_0\mu_0}$ midiendo la carga de una botella de Leyden
(primer capacitor), obteniendo $1/\sqrt{\varepsilon_0\mu_0}\approx
3\times10^8\ \text{m/s}$. Estas ondas electromagnéticas, cuya rapidez de
propagación es constante (pues $\varepsilon_0$ y $\mu_0$ son constantes
universales), fueron descubiertas experimentalmente por F.\,Hertz en la
Universidad de Karlsruhe en 1887. Otras medidas históricas de la rapidez
de la luz se deben a Galileo, O.\,Rømer (1675), J.\,Bradley (1727), y
Fizeau y Foucault (1849--1850), obteniendo $c=2.98\times10^8\ \text{m/s}$.
Estos resultados condujeron a Maxwell a postular que ``la luz es una
onda electromagnética''.
\end{observacion}

Volvamos a la pregunta: ¿satisfacen todas las leyes de la física el
principio de relatividad? La respuesta es: las ecuaciones de Maxwell
cambian su forma matemática bajo una transformación de Galileo, o sea,
\textbf{las ecuaciones de Maxwell no son invariantes de forma} bajo una
transformación de Galileo. Tampoco lo es la ecuación de ondas
electromagnéticas.

\begin{proposicionc}[No invariancia de la ecuación de onda]
La ecuación de onda electromagnética (13), en una dimensión espacial,
no es invariante de forma bajo una transformación de Galileo.
\end{proposicionc}

\begin{proof}
Partimos de $t'=t$, $x'=x+Vt$. Entonces
\[
\frac{\partial}{\partial t} = \frac{\partial t'}{\partial t}\frac{
\partial}{\partial t'} + \frac{\partial x'}{\partial t}\frac{\partial}
{\partial x'} = \frac{\partial}{\partial t'} + V\frac{\partial}{\partial
x'}, \qquad \frac{\partial}{\partial x} = \frac{\partial t'}{\partial
x}\frac{\partial}{\partial t'} + \frac{\partial x'}{\partial x}\frac{
\partial}{\partial x'} = \frac{\partial}{\partial x'}.
\]
Sustituyendo en $\partial^2\Psi/\partial x^2 - (1/c^2)\partial^2\Psi/
\partial t^2=0$ se obtiene, tras expandir,
\[
\left(1-\frac{V^2}{c^2}\right)\frac{\partial^2\Psi}{\partial x'^2} -
\frac{1}{c^2}\frac{\partial^2\Psi}{\partial t'^2} + \frac{2V}{c^2}
\frac{\partial^2\Psi}{\partial t'\partial x'} = 0,
\]
que es la ecuación de la onda en el sistema $K'$, la cual claramente no
tiene la misma forma que la ecuación original. \qedhere
\end{proof}

\subsection{La hipótesis del éter}

En resumen: las ecuaciones de Maxwell predicen la existencia de ondas
electromagnéticas. Para los físicos del siglo XIX debía existir un medio
de propagación de las ondas electromagnéticas, llamado \textbf{éter} (o
éter lumínifero, que significa ``portador de luz''). El S.R.I. con
respecto al cual las ecuaciones de Maxwell son particularmente simples
--el sistema elegido por Maxwell para presentar sus ecuaciones-- estaba
en reposo con respecto al éter, y fue llamado \textbf{sistema del éter}.
Las ecuaciones cambian a una forma más complicada si se refieren a un
S.R.I. que se traslada con velocidad constante con respecto al sistema
del éter; esto permitiría, en principio, determinar mediante
experimentos ópticos o eléctricos la rapidez de cualquier sistema con
respecto al éter.

Sea $K$ un S.R.I. fijo al éter. Con respecto a $K$, la rapidez de
propagación de la onda es $c$, y el movimiento de la onda con respecto
al éter es dado por
\[
\Psi(t,\vec x) = \Psi_0\cos\!\left(2\pi\nu\left(t-\frac{\hat n\cdot\vec
x}{c}\right)\right).
\]

\begin{proposicionc}
Consideremos los máximos de la onda, $\Psi=\Psi_0$, sobre el frente de
onda $\vec x_M(t)$. Derivando la condición $\Psi(t,\vec x_M)=\Psi_0$
respecto del tiempo se obtiene
\[
c = \vec V_M\cdot\hat n,
\]
donde $\vec V_M=d\vec x_M/dt$. El período de la onda es $T=1/\nu$ y su
longitud de onda es $\lambda=cT$.
\end{proposicionc}

\begin{proof}
De $\cos\big(2\pi\nu(t-\hat n\cdot\vec x_M/c)\big)=1$ se sigue
$2\pi\nu(t-\hat n\cdot\vec x_M/c)=2\pi n$, $n\in\mathbb{N}$. Derivando
respecto de $t$,
\[
2\pi\nu\left(1-\frac{d\vec x_M}{dt}\cdot\frac{\hat n}{c}\right)=0
\quad\Longrightarrow\quad 1-\frac{\vec V_M\cdot\hat n}{c}=0
\quad\Longrightarrow\quad c=\vec V_M\cdot\hat n . \qedhere
\]
\end{proof}

Sea ahora $K'$ un S.R.I. que se mueve con respecto al éter con rapidez
$V$ en la dirección del eje $x$ de $K$. En $K'$, la onda electromagnética
se escribe
\[
\Psi'(t',\vec x\,') = \Psi_0\cos\!\left(2\pi\nu'\left(t'-\frac{\hat
n\cdot\vec x\,'}{c'}\right)\right).
\]

\begin{proposicionc}[Efecto Doppler bajo Galileo]
Aplicando la transformación de Galileo $t'=t$, $\vec x = \vec x\,'+\vec
Vt'$ y comparando ambas expresiones de la onda, se obtiene
\[
\nu' = \nu\left(1-\frac{\hat n\cdot\vec V}{c}\right), \tag{14}
\]
donde $\nu$ es la frecuencia medida en el S.R.I. en reposo (respecto al
éter) y $\nu'$ es la frecuencia medida en el S.R.I. que se mueve con
velocidad $\vec V$. Además,
\[
c' = c - \hat n\cdot\vec V = c-V\cos\alpha . \tag{15}
\]
\end{proposicionc}

\begin{proof}
Sustituyendo $\vec x=\vec x\,'+\vec Vt'$ en $\Psi$ y comparando con
$\Psi'$,
\[
\nu'\left(t'-\frac{\hat n\cdot\vec x\,'}{c'}\right) = \nu\left(t'-
\frac{\hat n\cdot(\vec x\,'+\vec Vt')}{c}\right) \;\Longrightarrow\;
\nu't + \left(\frac{\nu}{c}-\frac{\nu'}{c'}\right)\hat n\cdot\vec x\,' =
\nu\left(1-\frac{\hat n\cdot\vec V}{c}\right)t .
\]
El número de onda es un absoluto (número de sucesos), por lo que
$1/\lambda=\nu/c=\nu'/c'$; sustituyendo se cancela el término en
$\hat n\cdot\vec x\,'$ y se obtiene (14). De $\nu/c=\nu'/c'$ se sigue
$\nu'=(c'/c)\nu$, que junto con (14) da (15). Las ecuaciones (14) y (15)
son válidas bajo la suposición de invariancia de Galileo. \qedhere
\end{proof}

\begin{observacion}[Fuente en movimiento]
Si la fuente se mueve con velocidad $\vec V_0$ respecto al éter, sea
$\nu_0$ su frecuencia propia (medida en su propio reposo) y $\nu$ la
medida respecto al S.R.I. del éter, de modo que
$\nu_0=\nu\big(1-\hat n\cdot\vec V_0/c\big)$. La frecuencia medida por un
S.R.I. que se mueve con velocidad $\vec V$ respecto al éter es
\[
\nu' = \nu\left(1-\frac{\hat n\cdot\vec V}{c}\right) = \nu_0\,
\frac{1-\hat n\cdot\vec V/c}{1-\hat n\cdot\vec V_0/c} .
\]
Con $\vec V_r=\vec V-\vec V_0$ la velocidad relativa de $K'$ respecto de
la fuente, y usando $V_r\ll c$,
\[
\nu' = \nu_0\left(1 - \frac{\hat n\cdot\vec V_r}{c} - \frac{(\hat
n\cdot\vec V_0)(\hat n\cdot\vec V_r)}{c^2} + \dots\right). \tag{16}
\]
La medición del término cuadrático posibilitaría la determinación de
$\vec V_0$, es decir, de la velocidad absoluta respecto del éter.
Johannes Stark llevó a cabo mediciones para comprobar (16) hasta el
término $\sim1/c$; para términos de orden superior sus medidas no fueron
precisas. El primer intento experimental para detectar efectos de
segundo orden en fenómenos ópticos debido al movimiento de la Tierra fue
llevado a cabo por Michelson en 1881, y más tarde refinado por Michelson
y Morley en 1887.
\end{observacion}

\section{Experimento de Michelson-Morley}

La idea es usar el interferómetro de Michelson para desdoblar un rayo de
luz de modo que los dos rayos resultantes viajen por caminos que forman
un ángulo recto uno con respecto al otro, para luego recombinarlos y
generar un patrón de franjas de interferencia, rotando el aparato en
$90^\circ$. Si la velocidad de la luz en el aparato depende de la
velocidad absoluta de la Tierra, entonces las velocidades de la luz a lo
largo de los dos caminos debieran cambiar cuando el instrumento es
rotado, y se debiera observar un corrimiento de las franjas.

\begin{center}
\begin{tikzpicture}[scale=0.85]
  \draw[thick] (0,0) -- (3.2,0) node[right] {$S_1$};
  \draw[thick] (0,0) -- (0,3.2) node[above] {$S_2$};
  \draw[fill=colordef!20,draw=colordef] (-0.15,-0.15) rectangle (0.15,0.15);
  \node at (0.3,-0.35) {$P$};
  \draw[thick,colorej] (0,-1.6) -- (0,-0.15) node[midway,right,black] {\footnotesize fuente};
  \draw[thick] (1.9,-1.6) -- (1.9,-0.4) node[above] {$T$};
  \node at (1.9,-1.9) {\footnotesize telescopio};
\end{tikzpicture}
\end{center}

El primer rayo se dirige hacia el espejo $S_1$ y luego, a través de $P$
(donde es reflejado), es recibido en el telescopio $T$. El segundo rayo
es reflejado primero en $S_2$ y luego se dirige a través de $P$ hacia el
telescopio $T$. Para lograr que el camino a través del espejo
semirreflectante sea igual para ambos rayos, es necesario introducir en
el camino del rayo hacia $S_2$ un vidrio compensador $G$.

Debe notarse que el aparato no puede construirse con la precisión
necesaria para que ambos rayos recorran exactamente el mismo camino; es
la \emph{diferencia} en la longitud de los dos rayos la que debe ser
observada en el telescopio $T$. Todo lo afirmado es válido si el
laboratorio está en reposo. Sin embargo: el laboratorio se encuentra en
la Tierra, la Tierra rota alrededor del Sol con rapidez de $30\
\text{km/s}$, y el Sol se mueve con rapidez de $300\ \text{km/s}$
alrededor de la constelación de Hércules. No podemos suponer que la
disposición experimental se encuentra en reposo con respecto al éter, y
un cambio en la orientación del dispositivo debiera conducir a un cambio
en el patrón de interferencia.

\subsection{Camino óptico a lo largo de la dirección del movimiento}

Consideremos primero el caso en que el aparato se mueve con velocidad
$\vec V$ en la dirección del rayo hacia $S_1$.

\begin{center}
\begin{tikzpicture}[scale=1]
  \draw[thick] (0,0) -- (4,0);
  \node at (-0.3,0) {$P$};
  \draw[thick] (3,-0.05) -- (3,0.05);
  \node at (3,0.35) {$S_1$};
  \draw[thick] (4,-0.05) -- (4,0.05);
  \node at (4,0.35) {$S_1'$};
  \draw[-Stealth,thick,colorej] (0,-0.6) -- (1.3,-0.6) node[midway,below] {$V$};
  \draw[decorate,decoration={brace,amplitude=4pt}] (0,0.5) -- (3,0.5)
    node[midway,above=2pt] {\footnotesize$l_1$};
  \draw[decorate,decoration={brace,amplitude=4pt}] (3,-1.2) -- (4,-1.2)
    node[midway,below=2pt] {\footnotesize$x$};
\end{tikzpicture}
\end{center}

Después de que la luz abandona el prisma $P$, el espejo $S_1$ recorre un
camino adicional $x$ antes de ser alcanzado por el rayo.

\begin{proposicionc}
El camino recorrido por la luz en la ida hacia $S_1$ es
\[
x = \frac{l_1\,V/c}{1-V/c},
\]
y en el camino de regreso, con $P$ recorriendo un camino $x'$ al
encuentro de la luz,
\[
x' = \frac{l_1\,V/c}{1+V/c}.
\]
\end{proposicionc}

\begin{proof}
El tiempo de vuelo de la luz es igual al tiempo de movimiento del espejo
$S_1$: $V=x/t_1 \Rightarrow t_1=x/V$. Por otro lado, el tiempo de vuelo
de la luz de $P$ a $S_1'$ es $t_1=(l_1+x)/c$. Igualando,
$x/V=(l_1+x)/c \Rightarrow x=l_1(V/c)/(1-V/c)$. Análogamente, el tiempo
de vuelo de $S_1$ a $P$ es $t_1'=x'/V$ y también $t_1'=(l_1-x')/c$, de
donde $x'=l_1(V/c)/(1+V/c)$. \qedhere
\end{proof}

\begin{proposicionc}
El camino óptico total recorrido por este rayo es
\[
S_1 = \frac{2l_1}{1-V^2/c^2} .
\]
\end{proposicionc}

\begin{proof}
\[
S_1 = l_1+x+l_1-x' = 2l_1 + \frac{l_1V/c}{1-V/c} - \frac{l_1V/c}{1+V/c}
= 2l_1 + l_1\, \frac{V}{c}\cdot\frac{2V/c}{1-V^2/c^2} = 2l_1\left(1 +
\frac{V^2/c^2}{1-V^2/c^2}\right) = \frac{2l_1}{1-V^2/c^2} . \qedhere
\]
\end{proof}

\subsection{Camino óptico perpendicular al movimiento}

Consideremos ahora el camino del rayo de luz reflejado en el espejo
$S_2$, que se mueve perpendicular a la dirección del rayo.

\begin{center}
\begin{tikzpicture}[scale=1]
  \draw[thick] (0,0) -- (0,3);
  \node at (0,3.3) {$S_2$};
  \draw[thick] (-0.05,0) -- (0.05,0);
  \node at (0.35,0) {$P$};
  \draw[thick,dashed] (1.1,0) -- (1.1,3);
  \node at (1.4,3.3) {$S_2'$};
  \draw[-Stealth,thick,colorej] (-0.6,0) -- (-0.6,1) node[midway,left] {$V$};
  \draw[thick,colorprop] (0,0) -- (1.1,3) node[midway,right] {\footnotesize$\sqrt{y^2+l_2^2}$};
  \draw[decorate,decoration={brace,amplitude=3pt,mirror}] (0,-0.3) -- (1.1,-0.3)
    node[midway,below=2pt] {\footnotesize$y$};
\end{tikzpicture}
\end{center}

Mientras la luz viaja de $P$ a $S_2$, el aparato recorre un camino $y$
en la dirección del movimiento; mientras la luz viaja de $S_2$ de vuelta
a $P$, el aparato recorre el mismo camino $y$.

\begin{proposicionc}
El camino óptico total recorrido por este rayo es
\[
S_2 = \frac{2l_2}{\sqrt{1-V^2/c^2}} .
\]
\end{proposicionc}

\begin{proof}
Sea $t_2$ el tiempo que demora la luz en ir de $P$ a $S_2'$, que también
es el tiempo que demora el aparato en recorrer la distancia $y$:
\[
t_2 = \frac{y}{V}, \qquad t_2 = \frac{\sqrt{y^2+l_2^2}}{c}
\quad\Longrightarrow\quad y = \frac{l_2V/c}{\sqrt{1-V^2/c^2}} .
\]
El camino óptico total es $S_2=2\sqrt{y^2+l_2^2}$; sustituyendo $y$ y
simplificando se obtiene $S_2=2l_2/\sqrt{1-V^2/c^2}$. \qedhere
\end{proof}

\subsection{La diferencia de camino óptico}

\begin{proposicionc}
La diferencia de camino recorrido por la luz, para la orientación
considerada, es
\[
\Delta S_a = S_1-S_2 = \frac{2l_1}{1-V^2/c^2} - \frac{2l_2}{\sqrt{1-
V^2/c^2}} \approx 2(l_1-l_2) + \big(2l_1-2l_2\big)\frac{V^2}{c^2} +
\dots, \qquad V\ll c.
\]
Esta cantidad no es simétrica en $l_1$ y $l_2$.
\end{proposicionc}

Consideremos ahora una rotación del dispositivo experimental en
$90^\circ$; la diferencia de longitud de camino cambiará. Sea $\Delta
S_b$ el cambio de camino en la nueva disposición; siguiendo el mismo
método,
\[
-\Delta S_b = 2(l_2-l_1) + \big(2l_2-2l_1\big)\frac{V^2}{c^2} + \dots
\]

\begin{proposicionc}
La diferencia de ambas diferencias de camino da la diferencia de camino
óptico entre las dos orientaciones del experimento,
\[
\Delta S = \Delta S_a - \Delta S_b = (l_1+l_2)\,\frac{V^2}{c^2} +
\dots,
\]
cantidad que \textbf{sí} es simétrica en $l_1$ y $l_2$. Este es un
efecto del movimiento del aparato --es decir, de la Tierra-- con
respecto al éter, y debiera conducir a un cambio observable en el patrón
de interferencia.
\end{proposicionc}

El cambio de diferencia de fase asociado es
\[
\Delta\phi = \frac{\nu}{c}\Delta S = \frac{\nu}{c}(l_1+l_2)\frac{V^2}
{c^2} + \dots,
\]
que debería ser observable para elecciones apropiadas de $\nu$, $l_1$ y
$l_2$. Sin embargo, \textbf{este efecto no pudo ser comprobado
experimentalmente por Michelson y Morley}. Esto significa que no se
pudo encontrar ninguna evidencia de la existencia del éter.

\subsection{La contracción de Lorentz-Fitzgerald}

En 1892--1893, Lorentz y Fitzgerald propusieron una explicación para los
resultados de Michelson-Morley.

\begin{definicion}[Hipótesis de contracción]
Las varillas métricas se contraen en la dirección del movimiento en un
factor $f=\sqrt{1-V^2/c^2}$,
\[
l_1' = l_1\sqrt{1-\frac{V^2}{c^2}}, \tag{17}
\]
y el tiempo se dilata en el mismo factor,
\[
\Delta t' = \frac{\Delta t}{\sqrt{1-V^2/c^2}} . \tag{18}
\]
\end{definicion}

\begin{proposicionc}
Bajo la hipótesis de contracción, $\Delta S' = \Delta S_a'-\Delta
S_b' = 0$.
\end{proposicionc}

\begin{proof}
Sustituyendo (17) en $S_1$,
\[
S_1' = \frac{2l_1'}{1-V^2/c^2} = \frac{2l_1\sqrt{1-V^2/c^2}}{1-V^2/c^2}
= \frac{2l_1}{\sqrt{1-V^2/c^2}},
\]
que ahora coincide en forma con $S_2$. Entonces
\[
\Delta S_a' = S_1'-S_2 = \frac{2l_1}{\sqrt{1-V^2/c^2}} - \frac{2l_2}
{\sqrt{1-V^2/c^2}} = \frac{2(l_1-l_2)}{\sqrt{1-V^2/c^2}},
\]
y, por el mismo argumento tras la rotación, $-\Delta S_b' =
2(l_2-l_1)/\sqrt{1-V^2/c^2}$. Por lo tanto
\[
\Delta S' = \Delta S_a'-\Delta S_b' = \frac{2(l_1-l_2)}{\sqrt{1-
V^2/c^2}} - \frac{2(l_1-l_2)}{\sqrt{1-V^2/c^2}} = 0 . \qedhere
\]
\end{proof}

De esta manera ningún efecto es observable y el experimento quedaría
aclarado: el concepto de éter es sorteado por medio de la hipótesis de
contracción de Lorentz-Fitzgerald. El experimento de Michelson-Morley
muestra que no es posible detectar el movimiento con respecto del éter.

\begin{observacion}
La hipótesis de la contracción \textbf{contradice} la transformación de
Galileo: contradice la hipótesis de que ``la distancia entre dos puntos
es absoluta'', y contradice la hipótesis de que ``el intervalo de
tiempo es absoluto''.
\end{observacion}

\subsection{La transformación de Lorentz}

Situación a fines del siglo XIX y comienzos del XX: las ecuaciones de la
óptica cambian su forma bajo una transformación de Galileo, por lo que,
por medio de fenómenos ópticos, sería posible detectar el movimiento de
un S.R.I. respecto de otro S.R.I. --en particular, respecto de un S.R.I.
fijo al éter. Sin embargo, el resultado del experimento de
Michelson-Morley fue negativo.

Se concluye que las ecuaciones de Maxwell deben ser invariantes de
forma bajo algún grupo de transformaciones de coordenadas que relacione
sistemas en movimiento relativo uniforme (M.R.U.).

\begin{definicion}[Transformación de Lorentz]
En 1904, Lorentz encontró un grupo de transformaciones que deja las
ecuaciones de Maxwell invariantes de forma:
\[
t' = \frac{1}{\sqrt{1-V^2/c^2}}\left(t+\frac{V}{c^2}x\right), \qquad
x' = \frac{1}{\sqrt{1-V^2/c^2}}(x+Vt), \qquad y'=y, \qquad z'=z,
\]
conocidas como \textbf{transformaciones de Lorentz}.
\end{definicion}

\begin{observacion}
Lorentz supuso que toda la materia es esencialmente de origen
electromagnético, y de allí dedujo la ley de contracción de las
varillas. Poincaré mostró que todas las ecuaciones de la electrodinámica
son invariantes de forma bajo transformaciones de Lorentz.

En 1905, Einstein formuló la \textbf{teoría de la relatividad especial},
obteniendo los resultados de Lorentz y Poincaré. En lugar de apoyarse en
la electrodinámica, Einstein demostró que son necesarios solo dos
postulados para aclarar el punto. En su desarrollo, Einstein reevaluó
las ideas de espacio y tiempo de la física clásica, y probó, por medio
de experimentos pensados, que la fuente de las limitaciones de la física
clásica estaba en el concepto de \textbf{simultaneidad}.
\end{observacion}

\subsection{El concepto de simultaneidad}

\begin{definicion}[Simultaneidad]
Dos sucesos son \textbf{simultáneos} cuando suceden o se desarrollan al
mismo tiempo.
\end{definicion}

El concepto de simultaneidad es claramente comprensible cuando decimos,
por ejemplo, que llegamos a la plaza en el mismo momento en que se
encendió el alumbrado; no parece ofrecer duda porque el suceso se
observa en el mismo lugar en que se produce. El concepto de
simultaneidad se usó para determinar la longitud geográfica en el mar
mediante relojes mecánicos de gran precisión, dando en el siglo XVIII
solución a un problema crucial para la navegación. La simultaneidad
está en acuerdo con el tiempo absoluto newtoniano, porque el tiempo es
absoluto y pasa de forma eterna y uniforme, con independencia de
cualquier suceso exterior.

La mecánica de Newton (1642--1727) decía que todos los sistemas
inerciales, que se desplazan entre sí con movimiento uniforme, eran
equivalentes para la descripción de los fenómenos mecánicos. Einstein
(1879--1955) extendió este axioma a todos los fenómenos físicos, lo que
implicaba la compatibilidad de esta física con la mecánica newtoniana.
Pero el principio relativista y la constancia de la velocidad de la luz
en todos los sistemas inerciales encerraban una dificultad y una
aparente contradicción. La simultaneidad de dos sucesos en la teoría de
la relatividad no se podía garantizar si los sucesos no se daban en el
mismo punto.

\begin{ejemplo}[El experimento mental del andén]
Supongamos que un observador $O$ se encuentra en el punto $M$ de un
andén rectilíneo, equidistante de dos farolas $A$ y $B$ situadas en el
andén, que pueden encenderse simultáneamente con un interruptor. Cuando
$O$ acciona el interruptor, la luz de las dos farolas llega
simultáneamente a $M$, puesto que la luz recorre el mismo espacio
($AM=MB$). Luego, para $O$, el encendido de las farolas es simultáneo.

Supongamos ahora que un segundo observador $O'$ circula por una vía
rectilínea paralela al andén con velocidad $\vec v$. Si $O$ acciona el
interruptor que enciende las farolas en el momento en que $O'$ pasa
frente al punto $M$ del andén, como el sistema de referencia de $O'$ se
mueve respecto al de $O$, y $O'$ avanza acercándose a la farola $B$, la
luz de ésta le llegará antes que la de $A$. Por lo tanto, \textbf{los
sucesos de encendido de las farolas no son simultáneos para $O'$}.
\end{ejemplo}

\begin{observacion}[Simultaneidad y causalidad]
La simultaneidad, desde la visión relativista, influye en el principio
de causalidad, ya que un suceso $F$ que causa el efecto $E$ puede
observarse como simultáneo con $E$ desde algún sistema de referencia, e
incluso, desde algún otro, podría observarse $E$ antes que $F$. Para
resolver este dilema hay que abandonar la idea de un tiempo absoluto: el
tiempo es relativo, una coordenada más.
\end{observacion}

\section{Postulados de Einstein}

En la teoría de Newton, el principio de relatividad establece: todos
los sistemas de referencia inerciales son equivalentes hasta donde los
experimentos mecánicos lo avalan (\textbf{principio de relatividad de
Galileo}); la mecánica es invariante de forma bajo transformaciones de
Galileo. Esto puede reformularse así: no existe ningún experimento
mecánico que permita determinar si un cuerpo está en reposo o en
movimiento uniforme y rectilíneo.

Einstein se dio cuenta de que el principio, tal como fue establecido,
es vacío, debido a que no existe tal cosa como un experimento
\emph{puramente} mecánico: aun a un nivel muy elemental, cualquier
experimento mecánico involucra una observación, y el acto de observar
es de carácter electromagnético. En realidad, cualquier experimento
involucra prácticamente todas las ramas de la física. Einstein dio el
paso lógico de remover la restricción de considerar solo experimentos
mecánicos:

\begin{definicion}[Postulado I: Principio de Relatividad de Einstein]
Todos los sistemas de referencia inerciales son equivalentes. Es decir,
la equivalencia entre los S.R.I. $K$ y $K'$ significa que las leyes de
la naturaleza se cumplen de la misma manera en $K$ como en $K'$. En
particular, no existe ningún experimento que permita determinar si un
cuerpo está en reposo o en movimiento uniforme y rectilíneo.
\end{definicion}

\begin{definicion}[Postulado II: Constancia de la velocidad de la luz]
La velocidad de la luz es la misma en todos los sistemas de referencia
inerciales.
\end{definicion}

\begin{observacion}
La mecánica de Newton afirma que si dos autos se desplazan, uno a
$100\ \text{km/h}$ y otro a $150\ \text{km/h}$, el segundo supera al
primero en $50\ \text{km/h}$. Por composición clásica de velocidades, la
luz del faro del segundo auto superaría la velocidad de la luz del
primero; pero el segundo postulado de Einstein dice que la velocidad de
la luz es constante en todos los sistemas inerciales y las velocidades
no se superponen. Por lo tanto, no la alcanzará --lo que entra en
contradicción con la mecánica newtoniana.
\end{observacion}

La constancia de la velocidad de la luz implica una nueva redefinición
de S.R.I.:

\begin{definicion}[Sistema de referencia inercial]
Un sistema de referencia inercial es aquel respecto del cual la luz se
propaga en el vacío en línea recta en cualquier dirección.
\end{definicion}

Si un sistema de referencia $K'$ gira o se mueve con movimiento
acelerado con respecto de un S.R.I. $K$, entonces con respecto de $K'$
la luz no se propaga en línea recta en todas las direcciones: con
respecto de $K$ la luz recorre una línea recta, pero con respecto de
$K'$ la luz recorre una curva. Si se acepta que el principio de
relatividad es válido para toda la física, entonces debemos abandonar
el Boost de Galileo. Se deben encontrar nuevas ecuaciones de
transformación que satisfagan los dos postulados de Einstein.

\subsection{El Boost de Lorentz}

El Boost de Lorentz puede obtenerse a partir de: el principio de
homogeneidad e isotropía del espacio, el postulado de relatividad, y el
principio de la constancia de la rapidez de la luz.

Sean $K$ y $K'$ dos S.R.I. tales que $K$ se mueve a lo largo del eje
$x'$ con respecto de $K'$ con rapidez $V$. (Usaremos letras mayúsculas
para la velocidad relativa entre S.R.I., y minúsculas para la velocidad
de un cuerpo respecto de algún S.R.I.)

En mecánica newtoniana, el Boost de Galileo es $t'=t$, $x'=x+Vt$,
$y'=y$, $z'=z$, donde el tiempo es absoluto. En relatividad especial el
tiempo es relativo y debe ser una coordenada, por lo que el Boost de
Lorentz, análogo al de Galileo, debe tener la forma general
\[
t' = \gamma(V)t + \alpha(V)x, \qquad x' = \beta(V)t+\lambda(V)x, \qquad
y'=y, \qquad z'=z, \tag{1}
\]
donde los coeficientes $\gamma(V),\alpha(V),\beta(V),\lambda(V)$ deben
depender de $V$ y $c$, tales que cuando $c\to\infty$ se recupere el
Boost de Galileo. El sistema debe ser equivalente si $K$ está en reposo
y $K'$ se mueve con rapidez $-V$, esto es
\[
t = \gamma(-V)t'+\alpha(-V)x', \qquad x=\beta(-V)t'+\lambda(-V)x',
\qquad y=y', \qquad z=z'. \tag{2}
\]
Las expresiones (1) y (2) son inversas.

\begin{proposicionc}[Composición relativista de velocidades, forma preliminar]
Sea una partícula que se mueve a lo largo del eje $x$ de $K$ y $x'$ de
$K'$, con velocidad $v=dx/dt$ en $K$ y $v'=dx'/dt'$ en $K'$. Usando (1),
\[
v' = \frac{\beta+\lambda v}{\gamma+\alpha v} . \tag{3}
\]
\end{proposicionc}

\begin{proof}
De (1), $dt'=\gamma\,dt+\alpha\,dx$, $dx'=\beta\,dt+\lambda\,dx$, de
donde $v'=dx'/dt' = (\beta\,dt+\lambda\,dx)/(\gamma\,dt+\alpha\,dx) =
(\beta+\lambda v)/(\gamma+\alpha v)$, dividiendo numerador y
denominador por $dt$. \qedhere
\end{proof}

Determinemos ahora los cuatro coeficientes.

\begin{proposicionc}
$\beta(V) = \lambda(V)\,V$, y $\lambda(V)=\gamma(V)$.
\end{proposicionc}

\begin{proof}
Si la partícula está en reposo respecto de $K'$, $v'=0$, entonces
respecto de $K$ se mueve con rapidez $-V$; sustituyendo $v=-V$, $v'=0$
en (3), $0=(\beta-\lambda V)/(\gamma-\alpha V) \Rightarrow
\beta=\lambda V$. Sustituyendo esto en (1)-(3), $x'=\lambda(Vt+x)$ y
$v'=\lambda(V+v)/(\gamma+\alpha v)$. Si en cambio la partícula se mueve
junto con $K$, $v=0$, respecto de $K'$ se mueve con velocidad $v'=V$;
sustituyendo, $V=\lambda V/\gamma \Rightarrow \lambda=\gamma$.
\qedhere
\end{proof}

Con esto, $t'=\gamma(V)t+\alpha(V)x$, $x'=\gamma(V)(Vt+x)$, y
$v'=\gamma(V+v)/(\gamma+\alpha v)$.

\begin{proposicionc}
$\alpha(V) = V/(c^2\gamma)$.
\end{proposicionc}

\begin{proof}
Si un haz de luz se mueve en $K$ con $v=c$, por el postulado de
constancia de la luz también $v'=c$; sustituyendo en la expresión
anterior, $c=\gamma(V+c)/(\gamma+\alpha c)$, de donde
$\alpha=V/(c^2\gamma)$. \qedhere
\end{proof}

Así,
\[
t' = \gamma(V)\left(t+\frac{V}{c^2}x\right), \qquad x' = \gamma(V)(Vt+x),
\qquad y'=y, \qquad z'=z, \qquad v' = \frac{V+v}{1+\dfrac{vV}{c^2}},
\]
donde la última ecuación es la \textbf{composición relativista de
velocidades}. Solo falta conocer $\gamma(V)$.

\begin{proposicionc}
\[
\gamma(V) = \frac{1}{\sqrt{1-V^2/c^2}} .
\]
\end{proposicionc}

\begin{proof}
Sustituyendo $t$ y $x$ (transformación inversa, con $-V$) en la
expresión de $t'$,
\[
t' = \gamma(V)\left[\gamma(-V)\left(t'-\frac{V}{c^2}x'\right) +
\frac{V}{c^2}\gamma(-V)(-Vt+x)\right] = \gamma(V)\gamma(-V)\left(1-
\frac{V^2}{c^2}\right)t',
\]
de donde
\[
1 = \gamma(V)\gamma(-V)\left(1-\frac{V^2}{c^2}\right). \tag{4}
\]
La equivalencia de los sistemas $K$ y $K'$, junto con el principio de
isotropía del espacio, nos dice que $\gamma$ no puede depender del
signo de la velocidad --si dependiera, un reloj fijo en $K'$
($\Delta x'=0$) mostraría $\Delta t=\gamma(-V)T'$ respecto de $K$, y si
$\gamma(-V)$ fuese negativo, $K$ observaría que el tiempo retrocede.
Isotropía implica entonces $\gamma=\gamma(V^2)$, y sustituyendo en (4)
se obtiene $\gamma(V^2) = 1/\sqrt{1-V^2/c^2}$. \qedhere
\end{proof}

\begin{definicion}[Transformación de Lorentz]
Las transformaciones de Lorentz son
\[
t' = \frac{t+\dfrac{V}{c^2}x}{\sqrt{1-V^2/c^2}}, \qquad x' = \frac{Vt+x}
{\sqrt{1-V^2/c^2}}, \qquad y'=y, \qquad z'=z,
\]
con inversa
\[
t = \frac{t'-\dfrac{V}{c^2}x'}{\sqrt{1-V^2/c^2}}, \qquad x =
\frac{-Vt'+x'}{\sqrt{1-V^2/c^2}}, \qquad y=y', \qquad z=z'.
\]
\end{definicion}

La constancia de la rapidez de la luz tiene como consecuencia la
existencia de una \textbf{rapidez máxima}, puesto que
$1-V^2/c^2\ge 0 \Rightarrow V\le c$: la velocidad máxima en la
naturaleza es la rapidez de la luz. Las transformaciones de Lorentz
implican que todos los sistemas inerciales están relacionados por
transformaciones de Lorentz con velocidad constante $V<c$.

\subsection{Contracción de longitud}

Consideremos una regla fija en $K$ con extremos $x_1$ y $x_2$. En $K'$,
sus extremos tienen coordenadas $x_1'$ y $x_2'$, relacionadas por
\[
x_1' = \frac{Vt_1+x_1}{\sqrt{1-V^2/c^2}}, \qquad x_2' = \frac{Vt_2+x_2}
{\sqrt{1-V^2/c^2}} .
\]

\begin{proposicionc}[Contracción de longitud]
Sea $l_0=x_2-x_1$ la longitud de reposo (medida en $K$, donde la regla
está fija) y $l=x_2'-x_1'$ la longitud medida en $K'$ (donde la regla se
mueve). Entonces
\[
l = l_0\sqrt{1-\frac{V^2}{c^2}} < l_0 .
\]
\end{proposicionc}

\begin{proof}
Para medir la longitud de la regla con respecto de $K'$ es necesario
conocer las coordenadas $x'$ de ambos extremos de modo simultáneo en
$K'$, es decir $t_1'=t_2'$. Igualando las expresiones anteriores con
$t_1'=t_2'$ se obtiene $t_1\ne t_2$ en $K$; sustituyendo y despejando,
\[
l = x_2'-x_1' = \frac{V(t_2-t_1)+(x_2-x_1)}{\sqrt{1-V^2/c^2}}
\]
y, usando la relación entre $t_2-t_1$ y $x_2-x_1$ que impone
$t_1'=t_2'$, se simplifica a $l=l_0\sqrt{1-V^2/c^2}$. \qedhere
\end{proof}

Esto significa que la longitud de un cuerpo en la dirección de su
movimiento, con velocidad constante $V$, es reducida por un factor
$\sqrt{1-V^2/c^2}$: es la \textbf{contracción de la longitud}. No hay
efecto de contracción en las direcciones transversales a la dirección
del movimiento.

\subsection{Dilatación del tiempo}

Sea un reloj fijo en $x=x_1$ en el sistema $K'$ que registra dos eventos
sucesivos, separados por un intervalo $\Delta t'=T$ en $K'$: los eventos
en $K'$ son $(x_1,t_1)$ y $(x_1,t_1+T)$.

\begin{proposicionc}[Dilatación del tiempo]
Sea $T'$ el intervalo de tiempo entre estos dos sucesos, medido desde
$K$ (donde el reloj se mueve con rapidez $V$). Entonces
\[
T' = \frac{T}{\sqrt{1-V^2/c^2}} > T .
\]
\end{proposicionc}

\begin{proof}
Por la transformación de Lorentz,
\[
t_1' = \frac{t_1+\frac{V}{c^2}x_1}{\sqrt{1-V^2/c^2}}, \qquad t_2' =
\frac{t_1+T+\frac{V}{c^2}x_1}{\sqrt{1-V^2/c^2}},
\]
de donde $T'=t_2'-t_1' = T/\sqrt{1-V^2/c^2}$. \qedhere
\end{proof}

Es decir, los relojes en movimiento marchan más lento en un factor
$(1-V^2/c^2)^{-1/2}$; este fenómeno se llama \textbf{dilatación del
tiempo}.

\begin{observacion}[Confirmaciones experimentales]
Esta predicción ha sido confirmada experimentalmente en múltiples
ocasiones a través de mediciones del efecto Doppler transversal
(experimento sugerido por Einstein en 1907): como consecuencia de la
dilatación temporal, la radiación emitida por un átomo en movimiento
transversal a la dirección de observación presenta una variación
relativa de la longitud de onda $\Delta\lambda/\lambda_0\approx
\gamma-1$. La primera confirmación se debe a Ives y Stilwell (1938),
usando átomos de hidrógeno ($v/c\approx0.005$, $\gamma-1\approx10^{-5}$).
En 2003, Saathoff \emph{et al.} usaron espectroscopía láser de iones
rápidos ($v/c\approx0.064$), y en 2007 Reinhardt \emph{et al.} mejoraron
la precisión con relojes ópticos atómicos. El récord actual lo tienen
Chou \emph{et al.}, quienes verificaron la dilatación temporal con
relojes atómicos de aluminio moviéndose a $v\approx10\ \text{m/s}$
($\gamma-1\approx4\times10^{-16}$). Otros métodos incluyen medir la vida
media de partículas en movimiento y comparar relojes atómicos en
movimiento relativo (``paradoja de los gemelos'').
\end{observacion}

\subsection{Composición de velocidades}

Consideremos que el sistema $K$ se mueve a lo largo del eje $x'$ con
velocidad constante $V$ respecto de $K'$. Sea $(v_x,v_y,v_z) =
(dx/dt,dy/dt,dz/dt)$ la velocidad de una partícula en $K$ (no
necesariamente a lo largo de un eje coordenado), y $(v_x',v_y',v_z')$ su
velocidad en $K'$.

\begin{proposicionc}[Composición relativista de velocidades]
\[
v_x' = \frac{V+v_x}{1+\dfrac{V}{c^2}v_x}, \qquad v_y' = \frac{v_y
\sqrt{1-V^2/c^2}}{1+\dfrac{V}{c^2}v_x}, \qquad v_z' = \frac{v_z\sqrt{1-
V^2/c^2}}{1+\dfrac{V}{c^2}v_x} .
\]
\end{proposicionc}

\begin{proof}
De la transformación de Lorentz, $dt' = (dt+\frac{V}{c^2}dx)/
\sqrt{1-V^2/c^2}$, $dx'=(Vdt+dx)/\sqrt{1-V^2/c^2}$, $dy'=dy$, $dz'=dz$.
Entonces
\[
v_x' = \frac{dx'}{dt'} = \frac{Vdt+dx}{dt+\frac{V}{c^2}dx} =
\frac{V+dx/dt}{1+\frac{V}{c^2}(dx/dt)} = \frac{V+v_x}{1+\frac{V}{c^2}
v_x},
\]
\[
v_y' = \frac{dy'}{dt'} = \frac{dy\sqrt{1-V^2/c^2}}{dt+\frac{V}{c^2}dx} =
\frac{(dy/dt)\sqrt{1-V^2/c^2}}{1+\frac{V}{c^2}(dx/dt)} = \frac{v_y
\sqrt{1-V^2/c^2}}{1+\frac{V}{c^2}v_x},
\]
y análogamente para $v_z'$. \qedhere
\end{proof}

Estas ecuaciones determinan la ley de composición (suma) de velocidades
en relatividad especial. En el límite $c\to\infty$ se recupera
$v_x'=v_x+V$, $v_y'=v_y$, $v_z'=v_z$, las ecuaciones de transformación de
velocidades de la teoría newtoniana.

\begin{ejemplo}
Si la partícula se mueve paralela al eje $x$ de $K$ ($v_y=v_z=0$),
entonces $v_x' = (V+v_x)/(1+Vv_x/c^2)$, $v_y'=v_z'=0$. Si además
$v_x=c$,
\[
v_x' = \frac{V+c}{1+\dfrac{Vc}{c^2}} = \frac{V+c}{1+\dfrac{V}{c}} =
\frac{c(V+c)}{c+V} = c,
\]
es decir, la rapidez de la luz es la misma en $K$ y en $K'$, confirmando
el segundo postulado.
\end{ejemplo}

\subsection{Propiedades de las transformaciones de Lorentz}

Escribimos las transformaciones de Lorentz en forma matricial. Definimos
las coordenadas
\[
x^0=ct, \quad x^1=x, \quad x^2=y, \quad x^3=z, \qquad x'^0=ct',\quad
x'^1=x',\quad x'^2=y',\quad x'^3=z',
\]
tal que la transformación de Lorentz se escribe
\[
x'^0 = \frac{x^0+\frac{V}{c}x^1}{\sqrt{1-V^2/c^2}}, \qquad x'^1 =
\frac{\frac{V}{c}x^0+x^1}{\sqrt{1-V^2/c^2}}, \qquad x'^2=x^2, \qquad
x'^3=x^3.
\]

\begin{definicion}[Matriz de Lorentz]
Se define la matriz de $4\times4$
\[
\Lambda^\mu_{\ \nu}(V) = \begin{pmatrix}
\dfrac{1}{\sqrt{1-V^2/c^2}} & \dfrac{V/c}{\sqrt{1-V^2/c^2}} & 0 & 0 \\[8pt]
\dfrac{V/c}{\sqrt{1-V^2/c^2}} & \dfrac{1}{\sqrt{1-V^2/c^2}} & 0 & 0 \\[8pt]
0 & 0 & 1 & 0 \\ 0 & 0 & 0 & 1
\end{pmatrix}, \qquad \mu,\nu=0,1,2,3,
\]
donde $\mu$ rotula las filas y $\nu$ las columnas, tal que la
transformación de Lorentz se escribe $x'^\mu = \Lambda^\mu_{\ \nu}(V)\,
x^\nu$.
\end{definicion}

\begin{proposicionc}
La transformación de Lorentz es invertible: $\det(\Lambda^\mu_{\
\nu}(V)) = 1$.
\end{proposicionc}

\begin{proof}
\[
\det(\Lambda^\mu_{\ \nu}(V)) = \frac{1}{1-V^2/c^2} - \frac{V^2/c^2}
{1-V^2/c^2} = \frac{1-V^2/c^2}{1-V^2/c^2} = 1 . \qedhere
\]
\end{proof}

La inversa es la misma transformación de Lorentz con velocidad $-V$:
$x^\mu = \Lambda^\mu_{\ \nu}(-V)x'^\nu$, de donde
$(\Lambda^{-1})^\mu_{\ \nu}(V) = \Lambda^\mu_{\ \nu}(-V)$.

\begin{proposicionc}[Grupo de Lorentz]
Las transformaciones de Lorentz forman un grupo, llamado \textbf{grupo
de Lorentz}, denotado $SO(1,3)$.
\end{proposicionc}

Las transformaciones de Lorentz se simplifican definiendo la
\textbf{rapidez reducida} $\beta:=V/c<1$ y el \textbf{factor de
contracción de Lorentz}
\[
\gamma(\beta) := \frac{1}{\sqrt{1-\beta^2}},
\]
tal que
\[
\Lambda^\mu_{\ \nu}(\beta) = \begin{pmatrix} \gamma & \gamma\beta & 0 &
0 \\ \gamma\beta & \gamma & 0 & 0 \\ 0&0&1&0 \\ 0&0&0&1 \end{pmatrix},
\qquad (\Lambda^{-1})^\mu_{\ \nu}(\beta) = \begin{pmatrix} \gamma &
-\gamma\beta & 0 & 0 \\ -\gamma\beta & \gamma & 0 & 0 \\ 0&0&1&0 \\
0&0&0&1 \end{pmatrix},
\]
de modo que $x'^0=\gamma(x^0+\beta x^1)$, $x'^1=\gamma(\beta x^0+x^1)$,
$x'^2=x^2$, $x'^3=x^3$.

\subsubsection{El ángulo de rapidez}

Definamos el parámetro adimensional $\phi$ (\textbf{ángulo de
rapidez}) mediante
\[
\tanh(\phi) = \beta < 1, \qquad \tanh(x) := \frac{\sinh x}{\cosh x} =
\frac{e^x-e^{-x}}{e^x+e^{-x}} < 1 .
\]
Entonces
\[
\gamma = \frac{1}{\sqrt{1-\beta^2}} = \frac{1}{\sqrt{1-\tanh^2\phi}} =
\cosh(\phi), \qquad \gamma\beta = \sinh(\phi),
\]
de modo que
\[
\Lambda^\mu_{\ \nu}(\phi) = \begin{pmatrix} \cosh\phi & \sinh\phi & 0 &
0 \\ \sinh\phi & \cosh\phi & 0 & 0 \\ 0&0&1&0 \\ 0&0&0&1 \end{pmatrix},
\qquad
\begin{aligned}
x'^0 &= \cosh(\phi)x^0+\sinh(\phi)x^1 \\
x'^1 &= \sinh(\phi)x^0+\cosh(\phi)x^1 \\
x'^2 &= x^2, \qquad x'^3=x^3.
\end{aligned} \tag{7}
\]

Para entender el significado geométrico de $\phi$, comparemos con una
rotación bidimensional. Sea $Oxyz$ y $O'x'y'z'$ rotado un ángulo
$\theta$ en torno al eje $z$:
\[
x' = x\cos\theta - y\sin\theta, \qquad y'=x\sin\theta+y\cos\theta,
\qquad z'=z. \tag{8}
\]

\begin{proposicionc}
La rotación (8) deja invariante $x^2+y^2+z^2$, la métrica de Euclides,
y forma un grupo llamado \textbf{grupo especial ortogonal}, $SO(3)$.
\end{proposicionc}

\begin{proof}
\[
x'^2+y'^2+z^2 = (x\cos\theta-y\sin\theta)^2+(x\sin\theta+y\cos\theta)^2
+z^2 = (\cos^2\theta+\sin^2\theta)(x^2+y^2)+z^2 = x^2+y^2+z^2 .
\]
Escribiendo la rotación en forma matricial, $x'^i=O^i_{\ j}x^j$, se
tiene $\det(O^i_{\ j})=\cos^2\theta+\sin^2\theta=1$ (invertible), y la
métrica de Euclides transforma como $g'_{ij} = (O^{-1})^k_{\
i}(O^{-1})^l_{\ j}\delta_{kl} = \delta_{ij}$, es decir, se preserva.
\qedhere
\end{proof}

\begin{proposicionc}[Rotación compleja]
Haciendo $\theta=i\phi$ en (8) y usando $\cos(i\phi)=\cosh(\phi)$,
$\sin(i\phi)=i\sinh(\phi)$, se recupera la transformación de Lorentz
(7): las transformaciones de Lorentz corresponden a una \textbf{rotación
compleja} con ángulo imaginario $\theta=i\phi$.
\end{proposicionc}

\begin{proof}
De $\cos x=(e^{ix}+e^{-ix})/2$ y $\sin x=(e^{ix}-e^{-ix})/(2i)$, con
$x=i\phi$,
\[
\cos(i\phi) = \frac{e^{-\phi}+e^{\phi}}{2} = \cosh(\phi), \qquad
\sin(i\phi) = \frac{e^{-\phi}-e^{\phi}}{2i} = i\sinh(\phi).
\]
Sustituyendo en (8), $x'=\cosh(\phi)x-i\sinh(\phi)y$,
$y'=i\sinh(\phi)x+\cosh(\phi)y$; multiplicando la primera por $i$,
$ix'=\cosh(\phi)(ix)+\sinh(\phi)y$, $y'=\sinh(\phi)(ix)+\cosh(\phi)y$,
que tiene la forma de (7) con $x^0\leftrightarrow ix$, $x^1
\leftrightarrow y$. \qedhere
\end{proof}

Es claro que la transformación de Lorentz no deja invariante la
cantidad $x^2+y^2+z^2$ (solo la deja invariante una rotación espacial).
¿Qué cantidad deja invariante la transformación de Lorentz? Como la
transformación involucra funciones hiperbólicas, probemos con
$(ct)^2-x^2-y^2-z^2$.

\begin{proposicionc}
La cantidad $(ct)^2-x^2-y^2-z^2$ es invariante bajo (7).
\end{proposicionc}

\begin{proof}
\[
(ct')^2-x'^2-y'^2-z'^2 = \big(\cosh\phi\,ct+\sinh\phi\,x\big)^2 -
\big(\sinh\phi\,ct+\cosh\phi\,x\big)^2 - y^2-z^2
\]
\[
= \big(\cosh^2\phi-\sinh^2\phi\big)(ct)^2 - \big(\cosh^2\phi-
\sinh^2\phi\big)x^2 - y^2-z^2 = (ct)^2-x^2-y^2-z^2 . \qedhere
\]
\end{proof}

Esta cantidad se escribe en forma matricial como
\[
(ct)^2-x^2-y^2-z^2 = X^t\eta X, \qquad \eta = \begin{pmatrix} 1&0&0&0\\
0&-1&0&0\\ 0&0&-1&0\\ 0&0&0&-1 \end{pmatrix},
\]
y esta forma se llama \textbf{forma hiperbólica}.

\section{Espacio de Minkowski}

La forma hiperbólica permite definir un nuevo tipo de elemento de línea,
que involucra el tiempo. Hay dos convenciones usuales para el elemento
de línea en relatividad especial:
\[
ds^2 = c^2dt^2-dx^2-dy^2-dz^2 \quad\text{(forma usada en física de
partículas)},
\]
\[
ds^2 = -c^2dt^2+dx^2+dy^2+dz^2 \quad\text{(forma usada en física
teórica)}.
\]
Ambas son equivalentes, difiriendo solo por un signo global. En este
curso usaremos la segunda, llamada \textbf{forma relativista}:
\[
ds^2 = -c^2dt^2+dx^2+dy^2+dz^2 .
\]

\begin{quote}
\emph{``Desde ahora, el espacio por sí mismo y el tiempo por sí mismo
están destinados a desvanecerse en meras sombras, y solo una especie de
unión de ambos preservará una realidad independiente.''}
\end{quote}

A esta unión se le llama \textbf{espacio-tiempo}. En el contexto
newtoniano, el tiempo está separado del espacio euclidiano
tridimensional: dado que $t'=t$, dos eventos simultáneos $P_1$ y $P_2$
tienen una distancia $\Delta s^2=(x_2-x_1)^2+(y_2-y_1)^2+(z_2-z_1)^2$
que es invariante bajo una transformación de Galileo, $\Delta
s'^2=\Delta s^2$.

En relatividad especial, un suceso queda definido por las tres
coordenadas espaciales y el tiempo en que ocurre. Minkowski introdujo un
espacio-tiempo cuadridimensional, llamado \textbf{espacio de Minkowski},
$M^4$, donde los sucesos se representan por puntos llamados
\textbf{puntos de universo}.

\subsection{El intervalo entre dos sucesos}

Consideremos dos S.R.I., $K$ y $K'$, y dos sucesos: $P_1$, la emisión de
una señal luminosa desde $(x_1,y_1,z_1)$ en el instante $t_1$ de $K$; y
$P_2$, la llegada de la señal a $(x_2,y_2,z_2)$ en el instante $t_2$ de
$K$. Como la señal se propaga a la rapidez de la luz, la distancia
recorrida es $\Delta s_t = c\Delta t = c(t_2-t_1)$, y también
$\Delta s_e = \sqrt{(\Delta x)^2+(\Delta y)^2+(\Delta z)^2}$. Igualando,
\[
c^2(\Delta t)^2 - (\Delta x)^2-(\Delta y)^2-(\Delta z)^2 = 0 .
\]
Observando los mismos sucesos desde $K'$, y dado que la velocidad de la
luz es la misma en ambos sistemas, se obtiene análogamente
$c^2(\Delta t')^2-(\Delta x')^2-(\Delta y')^2-(\Delta z')^2=0$. Es
decir, $\Delta s_t^2-\Delta s_e^2=0$ tanto en $K$ como en $K'$.

Si la velocidad de la señal no es la de la luz, $v<c$,
\[
\Delta s_t^2 - \Delta s_e^2 = c^2\Delta t^2 - v^2\Delta t^2 =
(c^2-v^2)\Delta t^2 > 0,
\]
es decir, $\Delta s_t^2\ne\Delta s_e^2$ en general.

\begin{definicion}[Intervalo entre dos sucesos]
Minkowski definió el \textbf{intervalo} entre dos sucesos como
\[
\Delta s^2 = c^2\Delta t^2 - \Delta x^2-\Delta y^2-\Delta z^2, \tag{9}
\]
(definición original), o equivalentemente, en la forma relativista que
usaremos,
\[
\Delta s^2 = -c^2\Delta t^2+\Delta x^2+\Delta y^2+\Delta z^2 . \tag{10}
\]
\end{definicion}

\begin{proposicionc}
El intervalo entre dos sucesos es invariante bajo transformaciones de
Lorentz.
\end{proposicionc}

\begin{proof}
Sean $K$ y $K'$ relacionados por $x'^0=\gamma(x^0+\beta x^1)$,
$x'^1=\gamma(\beta x^0+x^1)$, $x'^2=x^2$, $x'^3=x^3$. El intervalo
medido por $K'$ es
\[
\Delta s'^2 = -(\Delta x'^0)^2+(\Delta x'^1)^2+(\Delta x'^2)^2+(\Delta
x'^3)^2 .
\]
Sustituyendo,
\[
\Delta s'^2 = -\gamma^2(\Delta x^0+\beta\Delta x^1)^2 + \gamma^2(\beta
\Delta x^0+\Delta x^1)^2 + (\Delta x^2)^2+(\Delta x^3)^2
\]
\[
= -\gamma^2(1-\beta^2)(\Delta x^0)^2 + \gamma^2(1-\beta^2)(\Delta x^1)^2
+ (\Delta x^2)^2+(\Delta x^3)^2 = -(\Delta x^0)^2+(\Delta x^1)^2+
(\Delta x^2)^2+(\Delta x^3)^2 = \Delta s^2,
\]
usando $\gamma^2(1-\beta^2)=1$. Por lo tanto $\Delta s'^2=\Delta s^2$:
el intervalo entre dos sucesos es el mismo en todos los S.R.I. Esta
invariancia es la expresión matemática de la constancia de la velocidad
de la luz. \qedhere
\end{proof}

Esta definición está en acuerdo con la pérdida de la simultaneidad: si
dos eventos son simultáneos en $K$ ($\Delta t=0$), en $K'$ se tiene
\[
-(\Delta x'^0)^2+(\Delta x'^1)^2+(\Delta x'^2)^2+(\Delta x'^3)^2 =
(\Delta x^1)^2+(\Delta x^2)^2+(\Delta x^3)^2 \ne 0,
\]
en general, de donde $\Delta x'^0\ne 0$, es decir $\Delta t'\ne 0$: los
dos eventos no son simultáneos en $K'$.

Si los dos eventos son infinitesimalmente contiguos, el intervalo se
escribe
\[
ds^2 = -(dx^0)^2+(dx^1)^2+(dx^2)^2+(dx^3)^2. \tag{11}
\]
Este elemento de línea nos dice que estamos en un espacio
cuadridimensional. Denotamos las coordenadas de un punto de universo $P$
como $x^\mu=(x^0,x^1,x^2,x^3)$, $\mu=0,1,2,3$. El elemento de línea (11)
se escribe entonces
\[
ds^2 = \eta_{\mu\nu}\,dx^\mu dx^\nu, \tag{13}
\]
donde se ha introducido
\[
\eta_{\mu\nu} = \begin{pmatrix} -1&0&0&0\\0&1&0&0\\0&0&1&0\\0&0&0&1
\end{pmatrix}, \tag{14}
\]
llamada \textbf{métrica de Minkowski}. Esta métrica define el espacio
de Minkowski, con símbolos de Christoffel
\[
\Gamma^\lambda_{\mu\nu} = 0, \tag{15}
\]
puesto que $\eta_{\mu\nu}$ es constante. El signo menos en la componente
temporal se debe a la constancia de la velocidad de la luz --más
explícitamente, a la existencia de una velocidad máxima.

Una partícula moviéndose en el espacio cuadridimensional genera una
curva llamada \textbf{línea de mundo} (o línea de universo). La
pendiente de esta curva corresponde a la velocidad, que en cuatro
dimensiones tendrá cuatro componentes y se llamará
\textbf{cuadrivelocidad}. Como la velocidad de la luz es constante, la
velocidad de la partícula $\vec v$ --y por lo tanto la cuadrivelocidad--
no puede tomar cualquier valor: la magnitud de la cuadrivelocidad está
acotada, la pendiente de la línea de mundo está acotada, y las
partículas no pueden describir cualquier línea de mundo. Los tipos de
líneas de mundo se caracterizan entonces por la velocidad de las
partículas.

\section{Cono de Luz}

Volvamos al elemento de línea y escribamos las coordenadas en la forma
$x^\mu=(x^0,x^i)$, con $i=1,2,3$ las componentes espaciales, de modo que
la métrica de Minkowski se escribe compactamente
\[
\eta_{\mu\nu} = \begin{pmatrix} -1 & 0_{1\times3} \\ 0_{3\times1} &
\delta_{ij} \end{pmatrix}, \tag{16}
\]
y el elemento de línea se descompone como $ds^2=\eta_{\mu\nu}dx^\mu
dx^\nu = -c^2dt^2+\delta_{ij}dx^idx^j$.

\begin{proposicionc}
Sea una partícula que se mueve con velocidad $v^i$ en un S.R.I. $K$.
Entonces
\[
ds^2 = -c^2dt^2\left(1-\frac{v^2}{c^2}\right). \tag{17}
\]
\end{proposicionc}

\begin{proof}
\[
ds^2 = -c^2dt^2+\delta_{ij}dx^idx^j = -c^2dt^2\left(1-\frac{1}{c^2}
\frac{dx^i}{dt}\frac{dx^j}{dt}\delta_{ij}\right) = -c^2dt^2\left(1-
\frac{v^iv_i}{c^2}\right) = -c^2dt^2\left(1-\frac{v^2}{c^2}\right).
\qedhere
\]
\end{proof}

La expresión (17) permite clasificar los tipos de intervalo en
relatividad especial.

\begin{definicion}[Clasificación del intervalo]
Sean dos eventos $O$ y $P$, con $O$ en el origen.
\begin{itemize}
\item Si $v<c$, entonces $ds^2<0$: el intervalo es de \textbf{tipo
tiempo}. Es posible encontrar un S.R.I. \textbf{comóvil} respecto al
cual $O$ y $P$ tienen la misma posición espacial. No es posible
encontrar un S.R.I. respecto al cual $O$ y $P$ sean simultáneos.
\item Si $v=c$, entonces $ds^2=0$: el intervalo es de \textbf{tipo luz}.
No es posible encontrar un S.R.I. respecto al cual $O$ y $P$ sean
simultáneos, ni uno respecto al cual tengan la misma posición espacial.
Solo señales luminosas (o partículas que se muevan a la velocidad de la
luz) pueden conectar $O$ y $P$.
\item Si $v>c$, entonces $ds^2>0$: el intervalo es de \textbf{tipo
espacio}. Los eventos $O$ y $P$ no tienen conexión causal. Es posible
encontrar un S.R.I. respecto al cual $O$ y $P$ son simultáneos, e
incluso S.R.I. respecto a los cuales $P$ antecede a $O$, y viceversa: no
tienen un orden temporal absoluto. No es posible encontrar un S.R.I.
respecto al cual tengan la misma posición espacial.
\end{itemize}
\end{definicion}

\begin{observacion}
El carácter del intervalo es un concepto absoluto. Las partículas
físicas son aquellas con $v\le c$. El principio de causalidad de
Galileo-Newton establece que un evento no puede suceder antes que su
causa; este principio sigue siendo válido en relatividad especial: no
existe un S.R.I. donde la relación causal de dos eventos se invierta. Es
decir, si la causa de un evento sucede en $K$ en $(t_1,x_1)$ y tiene un
efecto en $(t_2,x_2)$ con $t_2>t_1$, cualquier transformación a $K'$
debe satisfacer $t_2'>t_1'$.
\end{observacion}

\begin{center}
\begin{tikzpicture}[scale=1.05]
  \fill[colordef!12] (0,0) -- (2.4,2.4) -- (-2.4,2.4) -- cycle;
  \fill[colordef!12] (0,0) -- (2.4,-2.4) -- (-2.4,-2.4) -- cycle;
  \fill[colorobs!10] (0,0) -- (2.4,2.4) -- (2.4,-2.4) -- cycle;
  \fill[colorobs!10] (0,0) -- (-2.4,2.4) -- (-2.4,-2.4) -- cycle;
  \draw[thick] (-2.4,-2.4) -- (2.4,2.4);
  \draw[thick] (2.4,-2.4) -- (-2.4,2.4);
  \draw[-Stealth] (-2.7,0) -- (2.7,0) node[right] {$x$};
  \draw[-Stealth] (0,-2.7) -- (0,2.7) node[above] {$ct$};
  \node at (0,-0.35) {\footnotesize$O$};
  \node[colordef] at (0,2.0) {\footnotesize Futuro absoluto};
  \node[colordef] at (0,-2.0) {\footnotesize Pasado absoluto};
  \node[colorobs] at (1.7,0) {\footnotesize Desconectado};
  \node[colorobs] at (-1.7,0) {\footnotesize Desconectado};
  \node at (2.5,2.5) {\footnotesize$a$};
  \node at (-2.5,2.5) {\footnotesize$b$};
  \node at (-2.5,-2.5) {\footnotesize$d$};
  \node at (2.5,-2.5) {\footnotesize$c$};
\end{tikzpicture}
\end{center}

Consideremos como origen un suceso $O$ cualquiera. Sobre las rectas $ab$
y $cd$ se tiene $x=\pm ct$; girando estas líneas se genera el
\textbf{cono de luz}, en cuyo interior $v=c$.
\begin{itemize}
\item En la región superior (aOc en el diagrama original), todos los
puntos satisfacen $-c^2t^2+x^2<0$ con $t>0$: el intervalo entre
cualquier suceso de esta región y $O$ es tipo tiempo; todos sus sucesos
ocurren \emph{después} de $O$ en todos los S.R.I., y la región se llama
\textbf{futuro absoluto} respecto de $O$.
\item En la región inferior ($bOd$), todos los puntos satisfacen
$-c^2t^2+x^2<0$ con $t<0$: todos sus sucesos ocurren \emph{antes} de
$O$ en todos los S.R.I., y la región se llama \textbf{pasado absoluto}
respecto de $O$.
\item En las regiones laterales ($dOa$ con $x<0$, $cOb$ con $x>0$),
todos los puntos satisfacen $-c^2t^2+x^2>0$: el intervalo es de tipo
espacio, $O$ no está causalmente conectado con ningún suceso de estas
regiones, y se llaman \textbf{absolutamente separadas} de $O$.
\end{itemize}

\subsection{Rotaciones espaciales y el grupo de Lorentz}

Las transformaciones de Lorentz contemplan tanto los boosts como las
rotaciones espaciales:
\[
\text{Boost de Lorentz:}\quad x'^\mu=\Lambda^\mu_{\ \nu}x^\nu, \quad
\Lambda^\mu_{\ \nu} = \begin{pmatrix}\gamma&\gamma\beta&0&0\\
\gamma\beta&\gamma&0&0\\0&0&1&0\\0&0&0&1\end{pmatrix}, \qquad
\text{Rotación espacial:}\quad \Lambda^\mu_{\ \nu} = \begin{pmatrix}
1&0\\0&O^i_{\ j}\end{pmatrix},
\]
donde $O^i_{\ j}$ es una matriz de rotación en torno a un eje. Entonces
\[
\big(\text{Transformaciones de Lorentz}\big) = \big(\text{Boost de
Lorentz}\big) \oplus \big(\text{Rotaciones espaciales}\big).
\]

\begin{proposicionc}
La transformación de Lorentz deja invariante la métrica de Minkowski:
$\Lambda^\alpha_{\ \mu}\,\eta_{\alpha\beta}\,\Lambda^\beta_{\
\nu} = \eta_{\mu\nu}$.
\end{proposicionc}

\begin{proof}
Como el elemento de línea es invariante, $ds'^2=\eta_{\alpha\beta}
dx'^\alpha dx'^\beta = \eta_{\alpha\beta}\Lambda^\alpha_{\ \mu}dx^\mu
\Lambda^\beta_{\ \nu}dx^\nu = \eta_{\mu\nu}dx^\mu dx^\nu = ds^2$, y como
esto vale para cualquier $dx^\mu$, se sigue $\Lambda^\alpha_{\ \mu}
\eta_{\alpha\beta}\Lambda^\beta_{\ \nu}=\eta_{\mu\nu}$. Verificación
explícita con el boost a lo largo de $x^1$:
\[
\Lambda^t\eta\Lambda = \begin{pmatrix}\gamma&\gamma\beta&0&0\\
\gamma\beta&\gamma&0&0\\0&0&1&0\\0&0&0&1\end{pmatrix}
\begin{pmatrix}-1&0&0&0\\0&1&0&0\\0&0&1&0\\0&0&0&1\end{pmatrix}
\begin{pmatrix}\gamma&\gamma\beta&0&0\\\gamma\beta&\gamma&0&0\\
0&0&1&0\\0&0&0&1\end{pmatrix} = \begin{pmatrix}-\gamma^2(1-\beta^2)&0&
0&0\\0&\gamma^2(1-\beta^2)&0&0\\0&0&1&0\\0&0&0&1\end{pmatrix} =
\eta, \qedhere
\]
usando $\gamma^2(1-\beta^2)=1$.
\end{proof}

Se dice que la transformación de Lorentz \textbf{preserva la métrica}.

\begin{definicion}[Espacio-tiempo de Minkowski]
Se define el \textbf{espacio-tiempo de Minkowski} $M^4$ como un espacio
cuadridimensional con coordenadas cartesianas $x^\mu$, $\mu=0,1,2,3$, y
métrica asociada
\[
\eta_{\mu\nu} = \begin{pmatrix}-1&0&0&0\\0&1&0&0\\0&0&1&0\\0&0&0&1
\end{pmatrix}, \qquad \eta^{\mu\nu} = \begin{pmatrix}-1&0&0&0\\0&1&0&0\\
0&0&1&0\\0&0&0&1\end{pmatrix} \quad\text{(métrica inversa)}.
\]
\end{definicion}

\begin{definicion}[Producto escalar y norma de cuadrivectores]
El producto escalar entre cuadrivectores $A^\mu$ y $B^\mu$ es
\[
A\cdot B := \eta_{\mu\nu}A^\mu B^\nu = -A^0B^0+A^iB^i,
\]
y la norma de $A^\mu$ es $\|A^\mu\| := \sqrt{A\cdot A} = \sqrt{-(A^0)^2+
\|A^i\|^2}$. Usaremos la notación $A^2 := A\cdot A = \eta_{\mu\nu}
A^\mu A^\nu$.
\end{definicion}

\begin{definicion}[Clasificación de cuadrivectores]
Un cuadrivector $A^\mu$ se dice:
\begin{itemize}
\item \textbf{tipo tiempo} si $A^2<0$ (dentro del cono de luz);
\item \textbf{tipo luz} si $A^2=0$ (sobre el cono de luz);
\item \textbf{tipo espacio} si $A^2>0$ (fuera del cono de luz).
\end{itemize}
\end{definicion}

\begin{ejemplo}
\begin{enumerate}[label=\alph*)]
\item $A^\mu=(1,0,0,0)$: $A^2=-1<0$, tipo tiempo.
\item $A^\mu=(1,0,1,0)$: $A^2=-1+1=0$, tipo luz.
\item $A^\mu=(1,0,1,\sqrt2)$: $A^2=-1+1+2=2>0$, tipo espacio.
\end{enumerate}
\end{ejemplo}

\subsection{Subir y bajar índices con la métrica de Minkowski}

La métrica $\eta_{\mu\nu}$ permite bajar índices: dado el cuadrivector
$A^\mu$ (tensor de orden $(1,0)$), se define $A_\mu:=\eta_{\mu\nu}A^\nu$.

\begin{proposicionc}
\[
A_0 = -A^0, \qquad A_i = A^i .
\]
\end{proposicionc}

\begin{proof}
$A_0=\eta_{0\nu}A^\nu=\eta_{00}A^0=-A^0$, y $A_i=\eta_{i\nu}A^\nu=
\eta_{ij}A^j=\delta_{ij}A^j=A^i$. \qedhere
\end{proof}

Análogamente, con la métrica inversa $\eta^{\mu\nu}$ se suben índices,
$A^\mu:=\eta^{\mu\nu}A_\nu$, obteniendo el mismo resultado: $A^0=-A_0$,
$A^i=A_i$.

\begin{observacion}
Para tensores de orden $(p,q)$ arbitrario, al subir o bajar un índice:
la componente temporal ($\mu=0$) cambia en un signo menos; la componente
espacial ($\mu=i$) no cambia de signo. Si en cambio se usara la métrica
de física de partículas, $\eta_{\mu\nu}=\text{diag}(1,-1,-1,-1)$, la
diferencia sería $A^0=A_0$, $A^i=-A_i$: se cambian \emph{tres}
cantidades de signo, en vez de solo una con la métrica relativista. En
lo que sigue usaremos siempre la métrica relativista.
\end{observacion}

\subsection{Boost de Lorentz en dirección arbitraria}

La transformación de Lorentz estudiada hasta aquí corresponde al caso en
que $K$ se mueve respecto de $K'$ a lo largo del eje $x'$. Generalicemos
al caso en que $K$ se mueve con velocidad $\vec V$ en una dirección
arbitraria $\hat n$, $\vec V=V\hat n$.

\begin{definicion}
La transformación de Lorentz generalizada es $x'^\mu=\Lambda^\mu_{\
\nu}(\vec V)x^\nu$, con
\begin{align*}
\Lambda^0_{\ 0}(\vec V) &= \frac{1}{\sqrt{1-V^2/c^2}}, \qquad
\Lambda^i_{\ 0}(\vec V) = \frac{V^i/c}{\sqrt{1-V^2/c^2}}, \\
\Lambda^0_{\ j}(\vec V) &= \frac{V_j/c}{\sqrt{1-V^2/c^2}}, \qquad
\Lambda^i_{\ j}(\vec V) = \delta^i_j + \left(\frac{1}{\sqrt{1-V^2/c^2}}
-1\right)\hat n^i\hat n_j .
\end{align*}
\end{definicion}

Con la notación $\beta^i:=V^i/c=\beta\hat n^i$ y
$\gamma=1/\sqrt{1-\beta^2}$, la transformación se reescribe
\[
\Lambda^0_{\ 0}=\gamma, \qquad \Lambda^i_{\ 0}=\gamma\beta^i, \qquad
\Lambda^0_{\ j}=\gamma\beta_j, \qquad \Lambda^i_{\ j} = \delta^i_j+
(\gamma-1)\hat n^i\hat n_j, \tag{19}
\]
o, en forma matricial,
\[
\Lambda^\mu_{\ \nu}(\beta^k) = \begin{pmatrix} \gamma & \gamma\beta_j
\\ \gamma\beta^i & \delta^i_j+(\gamma-1)\hat n^i\hat n_j \end{pmatrix},
\tag{20}
\]
con inversa $(\Lambda^{-1})^\mu_{\ \nu}(\beta^k) = \Lambda^\mu_{\
\nu}(-\beta^k)$ (se cambia el signo de los términos con un solo índice
$\beta$).

\begin{proposicionc}
Explícitamente,
\[
x'^0 = \gamma\big(x^0+\beta^ix_i\big), \qquad x'^i = \gamma\big(\beta^i
x^0+\hat n^i\hat n_jx^j\big) + x^i - \hat n^i\hat n_jx^j .
\]
\end{proposicionc}

\begin{proof}
De $x'^\mu=\Lambda^\mu_{\ \nu}(\beta^k)x^\nu$ con (19),
$x'^0=\Lambda^0_{\ 0}x^0+\Lambda^0_{\ i}x^i = \gamma x^0+\gamma\beta_i
x^i$, y $x'^i = \Lambda^i_{\ 0}x^0+\Lambda^i_{\ j}x^j = \gamma\beta^i
x^0 + (\delta^i_j+(\gamma-1)\hat n^i\hat n_j)x^j = \gamma\beta^ix^0+x^i+
(\gamma-1)\hat n^i\hat n_jx^j$, que se reordena en la forma enunciada.
\qedhere
\end{proof}

Esta forma no es muy iluminadora físicamente; reescribámosla de manera
más elegante. Sea $\hat n$ una dirección fija y $\vec A$ un vector
arbitrario; descomponemos $\vec A = \vec A_\parallel + \vec A_\perp$,
con $\vec A_\parallel\perp\vec A_\perp$ y $\vec A_\parallel=A_\parallel
\hat n$.

\begin{proposicionc}
\[
A_\parallel = \hat n\cdot\vec A, \qquad \vec A_\perp = \vec A - (\hat
n\cdot\vec A)\hat n = \hat n\times(\vec A\times\hat n) .
\]
\end{proposicionc}

\begin{proof}
Contrayendo $\vec A=A_\parallel\hat n+\vec A_\perp$ con $\hat n$, y
usando $\hat n\cdot\vec A_\perp=0$, $\hat n\cdot\hat n=1$, se obtiene
$\hat n\cdot\vec A = A_\parallel$. Despejando, $\vec A_\perp = \vec A -
(\hat n\cdot\vec A)\hat n$. En notación de índices,
$A^i_\perp = A^i-\hat n^j A_j\hat n^i = (\delta^i_j-\hat n^i\hat
n_j)A^j =: P^i_{\ j}A^j$, donde $P^i_{\ j}:=\delta^i_j-\hat n^i\hat n_j$
es el \textbf{operador de proyección} (satisface $P^i_{\ j}\hat n^j=0$ y
$P^2=P$). Usando la identidad de Levi-Civita
$\delta^i_j\delta^k_l-\delta^i_l\delta^k_j = \epsilon^{ikn}
\epsilon_{jln}$ y $\hat n^2=1$, se reescribe
\[
A^i_\perp = (\delta^i_j\hat n^k\hat n_k - \hat n^i\hat n_j)A^j =
\epsilon^{ikn}\epsilon_{jln}\hat n^k\hat n_lA^j = \epsilon^{iln}\hat
n_l\,\epsilon_{njk}A^j\hat n^k,
\]
que en notación vectorial es $\vec A_\perp = \hat n\times(\vec A\times
\hat n)$. \qedhere
\end{proof}

Por lo tanto, $\vec A = (\hat n\cdot\vec A)\hat n + \hat n\times(\vec
A\times\hat n)$.

\begin{proposicionc}[Boost de Lorentz en forma vectorial]
\[
x'^0 = \gamma\big(x^0+\beta x_\parallel\big), \qquad \vec x\,' =
\gamma\big(\beta x^0+x_\parallel\big)\hat n + \hat n\times(\vec
x\times\hat n) . \tag{31}
\]
\end{proposicionc}

\begin{proof}
De $x'^0=\gamma(x^0+\beta^ix_i)=\gamma(x^0+\beta\hat n^ix_i)=\gamma(x^0+
\beta x_\parallel)$, usando $\hat n^ix_i=x_\parallel$. Para $x'^i$,
usando $\beta^i=\beta\hat n^i$ y la identidad anterior con $\vec A=\vec
x$,
\[
x'^i = \gamma\beta x_\parallel'\hat n^i \cdot(\cdots) \;\to\;
x'^i = \gamma(\beta x^0+x_\parallel)\hat n^i + \big(\delta^i_j-\hat
n^i\hat n_j\big)x^j = \gamma(\beta x^0+x_\parallel)\hat n^i +
\epsilon^{iln}\hat n_l\epsilon_{njk}x^j\hat n^k,
\]
que en notación vectorial da (31). \qedhere
\end{proof}

\begin{observacion}
De (31) se observa que la componente de $x^\mu$ \textbf{paralela} a la
velocidad es afectada por la transformación de Lorentz,
\[
x'^0 = \gamma(x^0+\beta x_\parallel), \qquad x_\parallel' = \gamma
(\beta x^0+x_\parallel),
\]
mientras que la componente \textbf{perpendicular} a la velocidad
permanece invariante:
\[
\vec x\,'_\perp = \hat n\times(\vec x\times\hat n) = \vec x_\perp .
\]
Se puede rotar el sistema coordenado de modo que el eje $x$ coincida
con la dirección de la velocidad, $\hat n=\ihat$, recuperando la
transformación particular de Lorentz estudiada anteriormente.
\end{observacion}
